\family{Sapotaceae}{Shea family}
\label{fam:sapotaceae}

Trees or shrubs with white latex. Leaves alternate, simple and entire,
often crowded at the ends of the twigs, without stipules, often with
silky or rusty hairs beneath. Flowers small, in clusters in the leaf
axils or at the ends of the twigs. Fruit a berry with one to several
large, glossy seeds, each with a broad scar. Two species in the Arami
Uplands, both valued for their fruits.

\section*{Key to the species}
\begin{florakey}
	\lead{1a}{Leaves whitish beneath; fruit orange-yellow with 4 or 5
		glossy seeds; forest tree}{\taxon{Chrysophyllum albidum}}
	\lead{1b}{Leaves hairless beneath, crowded at the ends of stout twigs;
		bark thick and corky, cracked into blocks; fruit green to yellowish
		with usually one seed; savanna tree}{\taxon{Vitellaria paradoxa}}
\end{florakey}

\species{Chrysophyllum albidum}{G.Don}
\label{sp:SAP01}

\field{Local and trade names} Yoruba: Agbalumo; Igbo: Udara. Trade:
African star apple.

\field{Status} Native.

\field{Description} Evergreen tree to 35~m with a bole fluted at the
base. Bark grey-brown; slash with white latex. Leaves elliptic to
oblong, 8--20~cm long, whitish beneath. Flowers small, cream,
clustered in the leaf axils. Fruit a globose berry 3--5~cm across,
orange-yellow when ripe, with a sweet-sour pulp and 4 or 5 glossy
brown seeds.

\field{Phenology} Flowering March to May; fruiting December to March.

\field{Habitat} Lowland semi-deciduous forest and forest--savanna
mosaic; often kept or planted near houses.

\field{Uses} The fruit is eaten fresh and sold in large quantities in
markets during the dry season \citep{keay1989,orwa2009}.

\field{Specimens examined} Okeri (C4), 420~m, fr., 23~Dec 1934, \specimen{Harding}{1287}{OKF};
Okeri (C4), 420~m, fl., 11~Apr 1936, \specimen{Harding}{1423}{OKF};
Okeri (C4), 420~m, fl., 26~May 1936, \specimen{Harding}{1457}{OKF};
Aramiwo Foot (B2), 640~m, fr., 10~Mar 1963, \specimen{Olaniyan}{595}{SVN};
Ewe Stream gallery forest (B3), 580~m, fl., 18~May 1965, \specimen{Olaniyan}{719}{SVN};
Aramiwo Foot (B2), 640~m, fr., 7~Dec 2020, \specimen{Alade}{360}{AUH};
Okeri (C4), 420~m, fr., 14~Mar 2023, \specimen{Alade}{690}{AUH}.

\field{Distribution} Grid squares B2, B3, C4.

\begin{center}
\distmap{B/2, B/3, C/4}
\end{center}

\species{Vitellaria paradoxa}{C.F.Gaertn.}
\label{sp:SAP02}

\field{Local and trade names} Yoruba: Emi; Hausa: Kadanya. Trade:
Shea.

\field{Status} Native.

\field{Description} Deciduous tree to 15~m. Bark dark grey, thick and
corky, cracked into rectangular blocks. Twigs stout. Leaves oblong,
10--25~cm long, crowded at the ends of the twigs, with wavy margins.
Flowers creamy white, fragrant, clustered at the ends of the twigs,
appearing while the tree is leafless. Fruit an ellipsoid berry
4--5~cm long, green to yellowish, with a sweet pulp and usually one
large seed (the shea nut).

\field{Phenology} Flowering January to March; fruiting May to July.

\field{Habitat} Guinea savanna woodland; kept in farmland when land is
cleared, forming open parkland.

\field{Notes} This species has often been placed in the genus
\taxon{Butyrospermum}.

\field{Uses} Shea butter, pressed from the kernels, is used for
cooking and in cosmetics and is an important export. The fruit pulp is
eaten \citep{arbonnier2004,orwa2009}.

\field{Specimens examined} Asiri Crossing (A2), 740~m, fl., 15~Mar 1963, \specimen{Olaniyan}{599}{SVN};
Aramiwo Escarpment, summit (B1), 900~m, fr., 4~Jun 1963, \specimen{Olaniyan}{628}{SVN};
Okuta Hill (A4), 820~m, fr., 4~Jun 1964, \specimen{Olaniyan}{687}{SVN};
Aramiwo Escarpment, summit (B1), 900~m, fr., 7~Jul 1966, \specimen{Olaniyan}{826}{SVN};
Aramiwo Escarpment, summit (B1), 900~m, fr., 23~Jul 2019, \specimen{Alade}{226}{AUH};
Idofin (A3), 760~m, st., 15~Apr 2021, \specimen{Alade}{504}{AUH};
Asiri Crossing (A2), 740~m, fr., 17~Jul 2025, \specimen{Alade}{984}{AUH}.

\field{Distribution} Grid squares A2, A3, A4, B1.

\begin{center}
\distmap{A/2, A/3, A/4, B/1}
\end{center}
